\chapter{Conclusion}

\noindent Dormancy in microbial systems helps to maintain biodiversity and stabilise ecosystem processes \cite{JayT.Lennon2011Msbt}. As Arctic permafrost becomes permanently degraded,  frozen soil organic matter is being made available to decomposers, thus releasing millions of tons of greenhouse gases \cite{JahnMarkus2010Gcca}. Through application of the predictive model, this project is relevant to decision making on climate change mitigation and conservation prioritisation. \\

\noindent Bioprospecting for culturable Arctic bacteria with cold-active enzymes has played an important part in developing industrial, agricultural and medical processes \cite{srinivas2009bacterial}. Arctic microbes have useful agricultural applications in regulating optimal fish health \cite{HamiltonErinF2019AACM} and in food processing such as milk fermentation and ice-cream production \cite{cid2016properties}. Through this investigation there may be scope for understanding new and bio-technologically useful species. \\

\noindent The Arctic contains some of the highest fidelity Mars analogue sites in the world \cite{pollard2009overview} and provides us with exciting opportunities to develop our understanding of interplanetary biology. A previous study has found indications of microbial metabolism in Arctic permafrost brine lenses formed during the Quaternary \cite{GilichinskyD2003Swbw}. These supercooled echoniches serve as a model for extraterrestrial life on cryogenic planets. By studying dormancy mechanisms we may glean insight into the life processes of extinct or extant extraterrestrial organisms. \\

\noindent Throughout this project I will engage with professional bodies to disseminate our findings. I will apply for the Royal Society of Biology's (RSB) Outreach and Engagement grant scheme to share our work with colleagues and the public. Moreover, I will present our findings during Biology Week 2021 and 2022. I will apply for the RSB Big Biology Day grant scheme, so that I can take our research into schools to engage children, particularly those from underrepresented groups. I will apply for the Microbiology Society's Conference Grant to share our work with colleagues. Moreover, I will submit articles detailing our research to Microbiology Today and other appropriate publications.\\

\noindent In conclusion, this project has far reaching implications for climate research, biotechnology and as a terrestrial analogue for extraterrestrial environments. I will disseminate our work and findings through formal and informal routes, engaging both colleagues and the general public. 